\documentclass[a4paper,12pt]{article}
\usepackage{fontspec}\defaultfontfeatures{Ligatures=TeX}
% \usepackage{setspace}\setstretch{1.3} % \begin{spacing}{1.3}
\usepackage[a4paper,vmargin={3cm,3.5cm},hmargin={3cm,3cm}]{geometry}
%-----------------------------------------------------------------------------%
\usepackage{settings}
%-----------------------------------------------------------------------------%
%%% Title %%%
\title{Akaike Information Criterion: A Short Introduction%
\footnote{I assume the usual regularity conditions for maximum likelihood:
an interior and identifiable true parameter, three times differentiable log-densities,
nonsingular Fisher information, and sufficient moment conditions for the convergence of expectations below.
I also assume the reader is familiar with the \acrshort{ml} estimator, its asymptotic properties, and Fisher information.}}
\author{陳\,捷 Jesse C.\ Chen\\\texttt{\href{https://jessekelighine.com}{jessekelighine.com}}}
\date{\today}
%-----------------------------------------------------------------------------%

\begin{document}

\maketitle

\section{What's Wrong with \glsentrylong{ml}?}

\noindent
Let $Y\coloneqq(Y_{1},\ldots,Y_n)$ be a data set we observe and
let $f(\cdot\given\theta)$ be a probability density model where $\theta$ is the parameter.
Suppose that we try and fit model $f(\cdot\given\theta)$ with data $Y$ to obtain the estimate of the parameter $\hat\theta_Y$ defined as follows:
\begin{align}
	\hat\theta_{Y} \coloneqq \arg\max_{\theta} \log f(Y\given\theta). \tag{\acrshort{ml}}
\end{align}
What are we actually doing here?
We are supposing that \emph{if} $Y$ is drawn according to the density $f(\cdot\given\theta_{0})$ for some true parameter $\theta_0$,
\emph{then} $\hat\theta_{Y}$ is a good estimate for $\theta_{0}$.
This line of reasoning is extensively argued by Ronald Fisher, the inventor of the \gls{ml} method.

Yet, this approach has an obvious problem:
What if $Y$ follows another distribution with density function $g(\cdot\given\phi_{0})$?
We can, of course, also find the \gls{ml} estimate for $\phi_{0}$:
\begin{align*}
	\hat\phi_{Y} \coloneqq \arg\max_{\phi} \log g(Y\given\phi).
\end{align*}
In the spirit of \gls{ml},
a straightforward and naive approach is to see which model fits the data better,
which means comparing the fitted log-likelihoods
\begin{align}\label{eq:ml-compare}
	\log f(Y\given\hat\theta_{Y})
	\quad\text{and}\quad
	\log g(Y\given\hat\phi_{Y})
\end{align}
and choose the model with the higher log-likelihood.
This comparison is meaningful only if both models describe the same observations and
their densities use the same dominating measure.
This poses another problem:
since we only have one realization of the data $Y$,
we can always find some density function $h(\cdot\given\psi)$ that is carefully tailored to fit the data at hand $Y$ very well,
producing a high fitted likelihood $h(Y\given\hat\psi_{Y})$.
However, if another data set $X$ drawn from the same distribution as $Y$ is presented,
the model $h(\cdot\given\hat\psi_Y)$ may fit the new data $X$ poorly.
This is known as the problem of \textbf{overfitting}.

Luckily,
in describing \textbf{overfitting},
we are motivated to do \textbf{cross-validation},
i.e., to use another data set $X$ that is independent of $Y$ but follows the same distribution
to evaluate a parameter estimate obtained from $Y$.

\section{Deriving \texorpdfstring{\acrshort{aic}}{AIC}}

Let's switch back to $f(\cdot\given\theta)$, with fixed parameter dimension $k$ as $n\to\infty$.
For the derivation, let $X=(X_1,\ldots,X_n)$ and $Y=(Y_1,\ldots,Y_n)$ be independent
samples from $f(\cdot\given\theta_0)$, and write
\begin{align*}
	f(X\given\theta)\coloneqq\prod_{i=1}^n f(X_i\given\theta)
\end{align*}
for the joint density; define $f(Y\given\theta)$ similarly.
Instead of trying to estimate and compare the log-likelihoods in \eqref{eq:ml-compare},
we try to estimate their \textbf{cross-validated} versions
\begin{align*}
	\log f(X\given\hat\theta_{Y})
	\quad\text{and}\quad
	\log g(X\given\hat\phi_{Y}).
\end{align*}
We derive the correction for $f$; the same calculation applies to $g$ under analogous
correct-specification and regularity assumptions.
Thus, $\hat\theta_Y$ is fitted on $Y$ and evaluated on an independent data set $X$.
Since $X$ is unavailable, we instead approximate the expected \textbf{cross-validated}
log-likelihood $\E_{X,Y}[\log f(X\given\hat\theta_Y)]$.

First, expand $\log f(X\given\hat\theta_{Y})$ to second order around $\hat\theta_{X}$:
\begin{align*}
	\log f(X\given\hat\theta_{Y})
	&=
	\log f(X\given\hat\theta_{X}) \tag{zeroth order}\\
	&\mathrel{+}
	(\hat\theta_{Y}-\hat\theta_{X})\T
	\left[\frac{\partial\log f(X\given\hat\theta_{X})}{\partial\theta}\right] \tag{first order}\\
	&\mathrel{+}
	\frac12
	(\hat\theta_{Y}-\hat\theta_{X})\T
	\left[
	\frac{\partial^2\log f(X\given\hat\theta_{X})}{\partial\theta\partial\theta\T}
	\right]
	(\hat\theta_{Y}-\hat\theta_{X})
	+o_p(1). \tag{second order}
\end{align*}
The $o_p(1)$ remainder follows from the regularity conditions and
$\hat\theta_Y-\hat\theta_X=O_p(n^{-1/2})$.
The first-order term is exactly zero because $\hat\theta_{X}$ is an interior \acrshort{ml} estimator.
Thus,
\begin{gather}
	\log f(X\given\hat\theta_{Y})
	=
	\log f(X\given\hat\theta_{X})
	+
	\frac12
	(\hat\theta_{Y}-\hat\theta_{X})\T
	\JJ(\hat\theta_{X})
	(\hat\theta_{Y}-\hat\theta_{X})
	+o_p(1), \label{eq:taylor}\\
	\shortintertext{where}
	\JJ(\hat\theta_{X})
	\coloneqq
	\frac{\partial^2\log f(X\given\hat\theta_{X})}{\partial\theta\partial\theta\T}.
	\notag
\end{gather}
Here comes the key insight of \acrshort{aic}:
the cross-validated log-likelihood $\log f(X\given\hat\theta_{Y})$ is equal, up to $o_p(1)$, to
the fitted log-likelihood $\log f(X\given\hat\theta_{X})$ plus a correction term.
The correction term involves the Hessian $\JJ(\hat\theta_{X})$ and the difference between the two \acrshort{ml} estimators $\hat\theta_{Y}-\hat\theta_{X}$.
Now split the quadratic form into three parts:
\begin{align*}
	(\hat\theta_{Y}-\hat\theta_{X})\T
	\JJ(\hat\theta_{X})
	(\hat\theta_{Y}-\hat\theta_{X})
	&=
	(\hat\theta_{X}-\theta_{0})\T
	\JJ(\hat\theta_{X})
	(\hat\theta_{X}-\theta_{0}) \label{eq:part-a}\tag{a}\\
	&\mathrel+
	(\hat\theta_{Y}-\theta_{0})\T
	\JJ(\hat\theta_{X})
	(\hat\theta_{Y}-\theta_{0}) \label{eq:part-b}\tag{b}\\
	&\mathrel-
	2(\hat\theta_{Y}-\theta_{0})\T
	\JJ(\hat\theta_{X})
	(\hat\theta_{X}-\theta_{0}). \label{eq:part-c}\tag{c}
\end{align*}

We now take expectations.
This point is essential: part \eqref{eq:part-c} does \emph{not} generally converge to zero in probability.
Instead, the standard \acrshort{ml} limits give
\begin{align*}
	\begin{aligned}
		\sqrt n(\hat\theta_X-\theta_0)&\dto
		\mathcal N_k(0,\mathcal I(\theta_0)^{-1}),\\
		\sqrt n(\hat\theta_Y-\theta_0)&\dto
		\mathcal N_k(0,\mathcal I(\theta_0)^{-1}),
	\end{aligned}
	\quad\text{and}\quad
	\frac{\JJ(\hat\theta_X)}{n}&\pto-\mathcal I(\theta_0),
\end{align*}
where the first two limiting normal vectors are independent and
$\mathcal I(\theta_0)$ is the Fisher information for one observation.
Their independence and zero means imply, under the stated moment conditions,
\begin{align}\label{eq:cross-term}
	\E_{X,Y}\!\left[
	(\hat\theta_{Y}-\theta_{0})\T
	\JJ(\hat\theta_{X})
	(\hat\theta_{X}-\theta_{0})
	\right]\to0.
\end{align}

Parts \eqref{eq:part-a} and \eqref{eq:part-b} are similar in form:
\begin{align*}
	(\hat\theta_{Y}-\theta_{0})\T
	\JJ(\hat\theta_{X})
	(\hat\theta_{Y}-\theta_{0})
	&=
	\trace\left(
	n(\hat\theta_{Y}-\theta_{0})(\hat\theta_{Y}-\theta_{0})\T
	\frac{\JJ(\hat\theta_{X})}{n}
	\right),\\
	(\hat\theta_{X}-\theta_{0})\T
	\JJ(\hat\theta_{X})
	(\hat\theta_{X}-\theta_{0})
	&=
	\trace\left(
	n(\hat\theta_{X}-\theta_{0})(\hat\theta_{X}-\theta_{0})\T
	\frac{\JJ(\hat\theta_{X})}{n}
	\right).
\end{align*}
Asymptotic normality, the information equality, and convergence of the relevant moments therefore give
\begin{align}\label{eq:quadratic-terms}
	\E_{X,Y}[\text{part \eqref{eq:part-a}}]\to-k
	\quad\text{and}\quad
	\E_{X,Y}[\text{part \eqref{eq:part-b}}]\to-k.
\end{align}
Thus it is the \emph{expectation} of each part, rather than the random part itself,
that is asymptotically $-k$.
Taking expectations in \eqref{eq:taylor} and using
\eqref{eq:cross-term}--\eqref{eq:quadratic-terms}, we obtain
\begin{align}\label{eq:expected-cv}
	\E_{X,Y}[\log f(X\given\hat\theta_{Y})]
	=
	\E_X[\log f(X\given\hat\theta_{X})]-k+o(1).
\end{align}
Since $X$ and $Y$ have the same distribution,
$\log f(Y\given\hat\theta_Y)-k$ is therefore an asymptotically unbiased estimate
of the expected cross-validated log-likelihood.
This is the central correction behind \acrshort{aic}.
It is conventionally written as
\begin{align}
	\aic
	\coloneqq
	2k-2\log f(Y\given\hat\theta_{Y}). \label{eq:aic}\tag{\acrshort{aic}}
\end{align}
Its sign connects the criterion with information theory and \acrlong{kl}.
The simple $k$ correction uses correct specification and the information equality;
for a misspecified family, the analogous correction need not equal $k$.

\section{\texorpdfstring{\acrshort{aic}}{AIC}'s Connection with \texorpdfstring{\acrlong{kl}}{KL Divergence}}

\acrshort{kl} divergence is an information-theoretic measure of the discrepancy between two distributions.
For densities $p$ and $q$ with respect to the same dominating measure, it is defined as
\begin{align*}
	\KL(p\ggiven q)
	\coloneqq
	\int_{\mathcal X}\log\left[\frac{p(x)}{q(x)}\right]p(x)\dd{x},
\end{align*}
provided the integral is well defined.
The two main properties of \acrshort{kl} are
\begin{enumerate}
	\item $\KL(p\ggiven q)\geq0$.
	\item $\KL(p\ggiven q)=0$ if and only if $p=q$ almost everywhere.
\end{enumerate}
That is, $\KL(p\ggiven q)$ is small when $p$ and $q$ are similar.

In our case, we want to know the discrepancy between
the ``true'' joint density $f(\cdot\given\theta_{0})$ and
the fitted joint density $f(\cdot\given\hat\theta_{Y})$.
Because the latter depends on the random training sample $Y$, the relevant discrepancy is its expectation:
\begin{align*}
	&\E_Y\!\left[
	\KL(f(\cdot\given\theta_0)\ggiven f(\cdot\given\hat\theta_Y))
	\right]\\
	&\quad=
	\int_{\mathcal X}
	\log f(X\given\theta_0)f(X\given\theta_0)\dd{X}
	-
	\E_Y\!\left[
	\int_{\mathcal X}
	\log f(X\given\hat\theta_Y)f(X\given\theta_0)\dd{X}
	\right]\\
	&\quad=
	\text{constant}
	-\E_{X,Y}[\log f(X\given\hat\theta_Y)].
\end{align*}
Here $\mathcal X$ is the sample space of the full data set $X$.
The first integral is the negative entropy of the true distribution and is constant across candidate models.
The quantity $-\E_{X,Y}[\log f(X\given\hat\theta_Y)]$ is the expected \textbf{cross-entropy}.
Equation \eqref{eq:expected-cv} therefore shows that
\begin{align*}
	-2\log f(Y\given\hat\theta_Y)+2k
\end{align*}
is asymptotically unbiased for twice the expected cross-entropy.
Thus minimizing expected cross-entropy is equivalent to minimizing expected \acrshort{kl} divergence.

Measuring the discrepancy between $f(\cdot\given\theta_{0})$ and
$f(\cdot\given\hat\theta_{Y})$ makes intuitive sense:
the problem of \textbf{overfitting} can be understood as a large discrepancy between the ``true'' likelihood and the ``estimated'' likelihood.
In the original paper \parencite{akaike-1974}, \acrshort{aic} is motivated by \acrshort{kl}.
Hence, \acrshort{aic} is represented as the \emph{negative} of the corrected log-likelihood to match the sign of \textbf{cross-entropy}.
Thus, in practice, we select the model with the \emph{smallest} \acrshort{aic}.

\section{Why Times Two?}

If we consider a Gaussian model with $\theta=(\mu,\sigma^2)$,
the log-likelihood is written as
\begin{align*}
	\log f(X\given \theta)
	=
	- \frac{n}{2} \log(2\pi)
	- \frac{n}{2} \log\sigma^2
	- \frac{1}{2\sigma^2} \sum_{i=1}^{n} (X_{i} - \mu)^2.
\end{align*}
Multiplying the log-likelihood by two removes these factors of $\frac12$.
More generally, the factor of two is a conventional scale inherited from likelihood-ratio statistics and model deviance;
it does not affect the ranking of candidate models.
The minus sign turns maximization of corrected log-likelihood into minimization of \acrshort{aic}.
\asDemonstrated

\printbibliography

\end{document}
